\section{\large\textbf{Discussion}}
From Figure \ref{fig:porcentajeQ}, it can be observed that the average correct response rate does not correlate with the $Q$ factor. However, a notable difference between sources with directivities 2 and 4 is evident. This variation can be attributed to differences in the D/R parameter, as reported by \textcite{espectrorev}, which highlights a proportional relationship between both parameters. Additionally, similarities were found with previous studies on sound localization using third-octave filtered pink noise in reverberant rooms \parencite{rev1}.

For sources with directivities 1 and 2, however, the difference is contrary to expectations. This could be due to the exclusion of the source balloon, considering only the directivity factor $Q$, which could introduce a potential source of error. When examining this difference across the spectrum in Figure \ref{fig:porcentajeF}, notable variation is observed in the 400 and 500 Hz bands for the source with directivity 4. This suggests that, unlike other sources, localization ability improved considerably in these bands. In the 250 Hz band, the percentage of correct responses is close to 50\%, which does not strongly indicate correct localization. Above the 500 Hz band, there is minimal variation.

Moreover, results show minimal difference between directivities 1 and 2, suggesting that this increase in directivity does not significantly enhance localization ability within the frequency range investigated. Starting from the 800 Hz band, all directivities exhibit similar accuracy rates. This could be due to room reverberation, which decreases consistently, and, in higher third-octave bands, increases the $D/R$ ratio, according to \textcite{Hartmann}.

In terms of results segmented by demographic questions, a reduced localization ability was found for In-Ear headphones. Participants who took the test on a mobile phone also showed decreased performance. This may be due to the fact that all participants who completed the test on a mobile phone used In-Ear headphones, whereas only 3 participants used In-Ear headphones with a computer, suggesting linked results. This difference may be attributable to possible audio compression applied by mobile devices, either due to the platforms used or the device itself.
